% =============================================================================
\chapter{Kết luận}\label{chap:ketluan}
% =============================================================================

\section{Kết quả đạt được}

Nhóm đã xây dựng hoàn chỉnh chương trình quản lý kho hàng siêu thị bằng C++17, tự cài đặt mọi cấu trúc dữ liệu và
thuật toán sắp xếp, không dùng \code{vector}, \code{map}, \code{sort} của STL. Đối chiếu với yêu cầu của đề tài
(Bảng~\ref{tab:yeucau}):

\begin{table}[htbp]
\centering\small
\caption{Mức độ hoàn thành yêu cầu}
\label{tab:hoanthanh}
\begin{tabularx}{\linewidth}{@{}Y c@{}}
\toprule
\textbf{Yêu cầu} & \textbf{Kết quả} \\
\midrule
B-Tree bậc $m$: tìm, chèn (tách), xóa (mượn, gộp), duyệt; bất biến được kiểm tra sau mỗi thao tác & \ok \\
Hàng đợi FIFO các lô (DSLK); nhập tạo lô mới, xuất lấy từ lô cũ nhất và lấy tiếp các lô sau & \ok \\
Danh sách phiếu (DSLK); cập nhật tồn kho luôn khớp tổng các lô & \ok \\
Tìm theo tên (chuỗi con), liệt kê theo nhóm; Merge Sort tự cài & \ok \\
Cảnh báo tồn dưới ngưỡng; cảnh báo hết hạn trong 7 ngày sắp theo HSD & \ok \\
Báo cáo giá trị tồn theo nhóm; top 10 bán chạy trong tháng & \ok \\
Đọc tệp khi khởi động, ghi khi thoát; kiểm tra mã trùng, sai kiểu, sai miền & \ok \\
Dữ liệu mẫu 35 mặt hàng; bộ lớn $10^4$ và $10^5$; đo thời gian chạy & \ok \\
Không rò rỉ bộ nhớ (kiểm tra bằng AddressSanitizer/LeakSanitizer) & \ok \\
\emph{Nâng cao:} in cây sau mỗi lần tách/gộp nút & \ok \\
\emph{Nâng cao:} đếm số nút truy xuất với $m = 3, 5, 10$ trên $10^5$ mặt hàng & \ok \\
\bottomrule
\end{tabularx}
\end{table}

Thực nghiệm xác nhận lý thuyết: số nút truy xuất tăng theo $\log_t n$ và giảm khoảng 2,3 lần khi tăng bậc từ 3
lên 10 (13,25 xuống 5,84 nút với $10^5$ mặt hàng); chiều cao nằm trong cận lý thuyết. Chương trình được kiểm thử
bằng 7 bộ test tự động (hơn 40 nhóm ca bình thường, biên, sai) và kiểm thử ngẫu nhiên đối chiếu mô hình.

Qua đồ án, nhóm rút ra một số bài học:
\begin{itemize}
  \item \textbf{Chốt giao diện chung sớm.} Việc thống nhất \code{Types.h}, tên hàm và định dạng tệp từ ngày đầu
        giúp bốn mô-đun ghép lại được. Những chỗ chưa chốt kỹ (mã lô và mã phiếu trong tệp có tiền tố chữ trong
        khi chương trình dùng số nguyên; tên hàm ngày) đều phải sửa khi ghép.
  \item \textbf{Quy ước sở hữu bộ nhớ.} Quy định ``B-Tree chỉ giữ con trỏ, Inventory giải phóng'' tránh được lỗi
        giải phóng hai lần. Khi ghép, lời gọi hủy cây kèm xóa luôn mặt hàng (trái quy ước) bị trình biên dịch báo lỗi và
        được sửa lại để Inventory tự giải phóng; công cụ kiểm tra bộ nhớ xác nhận không còn rò rỉ.
  \item \textbf{Dữ liệu dẫn xuất không nên lưu.} Tồn kho tính lại từ các lô nên không bao giờ lệch. Kiểm tra tự
        động bộ dữ liệu mẫu đã phát hiện phiên bản đầu của bộ sinh dữ liệu ghi số lượng lô chưa trừ phần đã xuất.
\end{itemize}

\section{Hạn chế}

\begin{itemize}
  \item Toàn bộ dữ liệu nằm trong bộ nhớ; ưu điểm lớn nhất của B-Tree (ít lần đọc khối đĩa) chưa được khai thác.
  \item Tìm theo tên phải duyệt toàn bộ cây, $O(n)$ mỗi lần tìm.
  \item Xuất kho lấy đúng theo thứ tự nhập (FIFO) mà không bỏ qua lô đã quá hạn; cảnh báo HSD chỉ thông báo cho
        người quản lý.
  \item Mã vạch chỉ được kiểm tra độ dài và chữ số, chưa kiểm tra chữ số kiểm tra (check digit) của chuẩn EAN.
  \item Khi đọc tệp, trùng mã lô giữa các mặt hàng khác nhau chưa được phát hiện (cần thêm một chỉ mục mã lô).
  \item Giao diện dòng lệnh, tiếng Việt không dấu; một người dùng tại một thời điểm.
\end{itemize}

\section{Hướng phát triển}

\begin{itemize}
  \item Lưu B-Tree trên đĩa, mỗi nút là một khối, chỉ nạp các nút cần thiết; dùng B$^+$-Tree (dữ liệu ở lá, các lá
        nối với nhau) để duyệt theo khoảng mã vạch nhanh hơn.
  \item Thêm chỉ mục cho tên hàng (cây tiền tố hoặc chỉ mục ngược theo từ) để tìm theo tên nhanh hơn $O(n)$.
  \item Xuất kho theo nguyên tắc FEFO (hết hạn trước, xuất trước) và tự động loại lô đã quá hạn.
  \item Kiểm tra chữ số kiểm tra EAN; giao diện đồ họa; nhiều người dùng.
\end{itemize}

% =============================================================================
% Tai lieu tham khao la mot muc cua Chuong 6 (theo de cuong), khong mo trang moi
\begingroup
\renewcommand{\chapter}[2]{\section*{#2}}
\makeatletter\let\@mkboth\@gobbletwo\makeatother   % giu tieu de dau trang cua Chuong 6
\small
\begin{thebibliography}{9}
\addcontentsline{toc}{section}{Tài liệu tham khảo}
\setlength{\itemsep}{1pt}

\bibitem{clrs} T. H. Cormen, C. E. Leiserson, R. L. Rivest, C. Stein.
  \emph{Introduction to Algorithms}, 4th ed. MIT Press, 2022. Chương 18: B-Trees.

\bibitem{bayer} R. Bayer, E. McCreight. Organization and Maintenance of Large Ordered Indexes.
  \emph{Acta Informatica}, 1(3):173--189, 1972.

\bibitem{comer} D. Comer. The Ubiquitous B-Tree. \emph{ACM Computing Surveys}, 11(2):121--137, 1979.

\bibitem{knuth} D. E. Knuth. \emph{The Art of Computer Programming, Vol. 3: Sorting and Searching}, 2nd ed.
  Addison-Wesley, 1998.

\bibitem{avl} G. M. Adelson-Velsky, E. M. Landis. An algorithm for the organization of information.
  \emph{Soviet Mathematics Doklady}, 3:1259--1263, 1962.

\bibitem{yao} A. C.-C. Yao. On random 2--3 trees. \emph{Acta Informatica}, 9(2):159--170, 1978.

\bibitem{sedgewick} R. Sedgewick, K. Wayne. \emph{Algorithms}, 4th ed. Addison-Wesley, 2011.

\bibitem{xorshift} G. Marsaglia. Xorshift RNGs. \emph{Journal of Statistical Software}, 8(14):1--6, 2003.

\end{thebibliography}
\endgroup
